\section{Introduction}
\label{sec:introduction}

Ocular artifacts are difficult to remove from electroencephalography (EEG)
because their spatial propagation changes with montage, reference, electrode
contact, and user physiology.  A population denoiser can learn common artifact
statistics, but a target user may have a different EOG-to-EEG transfer.  The
central question in this paper is deliberately narrow: can a short,
query-disjoint calibration recording define an unseen user's corruption
coordinates, and does iterative diffusion add value over a one-step estimator
given exactly the same coordinates and targets?

We follow the principle \emph{personalize the corruption, not the brain}.  The
clean-EEG model is never conditioned on participant identity.  Instead, a
rank-two artifact basis is estimated from support EEG/EOG by ridge regression,
thin SVD, and within-montage Procrustes alignment.  This basis has a structural
role: it defines the low-dimensional diffusion state, the observation supplied
at every reverse step, and the final EEG correction.  Query EOG, eye tracking,
artifact labels, outcomes, and participant identifiers are excluded from
inference.

The resulting support-calibrated artifact-subspace diffusion model predicts
bounded ocular coefficients rather than regenerating the full EEG.  We compare
it with an information-matched deterministic network trained on the identical
target and with population and same-cell wrong-operator interventions using a
shared checkpoint and common random numbers.  This design separates two claims:
(i) an incremental value of iterative diffusion, and (ii) an incremental value
of matching-user calibration.

Experiments use paired real-EEG/EOG-backed semi-simulation for waveform fidelity
and natural SGEYESUB recordings for participant-level artifact and preservation
proxies.  The evidence is intentionally reported as development evidence.  The
current experiment finds a clear advantage of the diffusion point estimate over
the matched one-step estimator on the paired task, together with useful
posterior risk ranking and acceptable natural-EEG safety.  However, matching
support does not reliably improve over population or same-cell wrong support.
Thus the current data support artifact-subspace diffusion, but do not establish
the stronger subject-aware claim.

Our contributions are:
\begin{itemize}
  \item a query-disjoint corruption calibration whose basis directly defines
  diffusion coordinates and reconstruction, without participant embeddings;
  \item an information-matched deterministic comparison and operator
  interventions that isolate diffusion utility from subject-calibration value;
  \item participant/source-record paired evaluation with three training seeds,
  posterior uncertainty analysis, and explicit natural-EEG preservation tests;
  \item a scoped negative finding: under the tested support estimator, the
  matching-user increment is not identified even though iterative diffusion
  improves over the one-step estimator.
\end{itemize}
